\subsection{对抗式交替优化训练策略}
\label{sec:adv_alternating_training_cn}

设一个训练批次记为 $\mathcal{B}=\{\mathbf{x}_i\}_{i=1}^{B}$，其中
$\mathbf{x}_i \in \mathbb{R}^{N \times C}$ 表示第 $i$ 个点云样本，$N$ 为点数，$C$ 为点特征维度。记位姿识别描述符网络为
$f_{\theta}(\cdot)$，其参数为 $\theta$；记动态遮挡生成器为 $G_{\phi}(\cdot)$，其参数为 $\phi$。
生成器根据输入点云预测一组具有车辆形态先验的三维长方体遮挡体，并通过几何遮挡机制对原始点云进行结构化扰动，从而构造用于鲁棒训练的对抗样本。

从优化角度看，该问题可表述为如下极小极大问题：
\begin{equation}
\min_{\theta}\max_{\phi}\;
\mathcal{L}_{\mathrm{place}}\!\left(f_{\theta}\!\left(T_{\phi}(\mathbf{x})\right)\right),
\label{eq:cn_minmax_objective}
\end{equation}
其中，$T_{\phi}(\cdot)$ 表示由生成器参数化的遮挡变换。考虑到生成器需要满足基本几何合理性约束，而描述符网络同时需要兼顾干净样本性能与受扰样本鲁棒性，本文采用交替优化方式分别更新 $G_{\phi}$ 与 $f_{\theta}$。

\paragraph{1) 批次内困难三元组位置识别损失}
对于批次中的任意锚样本 $i$，记其正样本集合与负样本集合分别为 $\mathcal{P}(i)$ 和 $\mathcal{N}(i)$。给定描述符
$\mathbf{z}_i=f_{\theta}(\mathbf{x}_i)$，定义样本间欧氏距离为
\begin{equation}
d_{ij}=\left\lVert \mathbf{z}_i-\mathbf{z}_j \right\rVert_2.
\label{eq:cn_pair_distance}
\end{equation}
采用批次内困难三元组损失（batch-hard triplet loss），即对每个锚样本选取批次内最难正样本和最难负样本，得到
\begin{equation}
\mathcal{L}_{\mathrm{place}}
=
\frac{1}{|\mathcal{V}|}
\sum_{i \in \mathcal{V}}
\left[
\max_{p \in \mathcal{P}(i)} d_{ip}
-
\min_{n \in \mathcal{N}(i)} d_{in}
+
\delta
\right]_+,
\label{eq:cn_batch_hard_triplet}
\end{equation}
其中，$\delta$ 为间隔超参数，$[\cdot]_+=\max(\cdot,0)$，$\mathcal{V}$ 表示当前批次中同时具有至少一个正样本和至少一个负样本的有效锚样本集合。

\paragraph{2) 遮挡生成器的前向建模}
对输入点云 $\mathbf{x}_i$，生成器预测 $M$ 个候选遮挡体：
\begin{equation}
G_{\phi}(\mathbf{x}_i)
\rightarrow
\left\{
\left(\mathbf{c}_{i,m}, \mathbf{s}_{i,m}, \psi_{i,m}\right)
\right\}_{m=1}^{M},
\label{eq:cn_generator_box_param}
\end{equation}
其中，$\mathbf{c}_{i,m}\in\mathbb{R}^{3}$、$\mathbf{s}_{i,m}\in\mathbb{R}^{3}$ 与 $\psi_{i,m}$ 分别表示第 $m$ 个长方体的中心、尺寸与偏航角。
在训练过程中，对每个样本随机采样一个激活遮挡体数量 $k_i \in \{1,\dots,M\}$，并在该样本上激活恰好 $k_i$ 个遮挡体，以形成不同强度的结构化扰动。

进一步地，生成器同时输出软遮挡掩码与硬遮挡掩码，因此可以定义两种扰动形式：
\begin{equation}
\tilde{\mathbf{x}}^{\,s}_i=T^{\mathrm{soft}}_{\phi}(\mathbf{x}_i), \qquad
\tilde{\mathbf{x}}^{\,h}_i=T^{\mathrm{hard}}_{\phi}(\mathbf{x}_i).
\label{eq:cn_soft_hard_transform}
\end{equation}
其中，$T^{\mathrm{soft}}_{\phi}$ 用于生成器更新阶段，以保证梯度能够从描述符损失反向传播到生成器；而 $T^{\mathrm{hard}}_{\phi}$ 用于描述符网络更新阶段，以使网络在更接近真实离散遮挡结果的样本上进行训练。若启用物体插入机制，则在硬遮挡变换阶段还会在激活长方体表面采样并插入合成点，以模拟动态障碍物的激光回波。

\paragraph{3) 生成器损失函数}
在生成器更新阶段，固定描述符网络参数 $\theta$，仅优化生成器参数 $\phi$。此时生成器的目标是在几何先验约束下，尽可能增大描述符网络在受扰点云上的位置识别损失。于是生成器损失定义为
\begin{equation}
\mathcal{L}_{G}
=
-\mathcal{L}_{\mathrm{place}}
\left(
\left\{
f_{\theta}\!\left(\tilde{\mathbf{x}}^{\,s}_i\right)
\right\}_{i \in \mathcal{B}}
\right)
+
\lambda_{\mathrm{size}} \mathcal{L}_{\mathrm{size}}
+
\lambda_{\mathrm{height}} \mathcal{L}_{\mathrm{height}}
+
\lambda_{\mathrm{range}} \mathcal{L}_{\mathrm{range}}.
\label{eq:cn_generator_loss}
\end{equation}
式中，第一项前的负号表示在最小化 $\mathcal{L}_{G}$ 的同时，生成器实际上在最大化对抗样本上的位置识别损失，从而学习更具破坏性的遮挡模式。后三项为几何正则约束，其具体定义如下。

若长方体尺寸可学习，则采用尺寸先验约束其接近典型车辆尺寸 $\mathbf{s}_{\mathrm{car}}$：
\begin{equation}
\mathcal{L}_{\mathrm{size}}
=
\begin{cases}
\mathrm{SmoothL1}\!\left(\mathbf{s}_{i,m}, \mathbf{s}_{\mathrm{car}}\right), & \text{长方体尺寸可学习时}, \\
0, & \text{采用固定尺寸时}.
\end{cases}
\label{eq:cn_size_prior}
\end{equation}
高度先验用于约束遮挡体中心靠近路面附近的合理垂直区域：
\begin{equation}
\mathcal{L}_{\mathrm{height}}
=
\frac{1}{BM}
\sum_{i=1}^{B}\sum_{m=1}^{M}
\left(c_{i,m}^{(z)}+1.0\right)^2.
\label{eq:cn_height_prior}
\end{equation}
距离先验用于抑制生成器将遮挡体放置在异常远距离区域：
\begin{equation}
\mathcal{L}_{\mathrm{range}}
=
\frac{1}{BM}
\sum_{i=1}^{B}\sum_{m=1}^{M}
\left[
\max\!\left(
0,\sqrt{\left(c_{i,m}^{(x)}\right)^2+\left(c_{i,m}^{(y)}\right)^2}-45.0
\right)
\right]^2.
\label{eq:cn_range_prior}
\end{equation}
因此，生成器的优化目标可以概括为：在满足合理几何分布的前提下，尽可能产生对描述符网络最不利的结构化动态遮挡。

\paragraph{4) 描述符网络损失函数}
在描述符更新阶段，固定生成器参数 $\phi$，重新利用同一遮挡预算为每个样本生成硬遮挡点云 $\tilde{\mathbf{x}}^{\,h}_i$，并分别计算干净样本与受扰样本的描述符：
\begin{equation}
\mathbf{z}^{\,c}_i=f_{\theta}(\mathbf{x}_i), \qquad
\mathbf{z}^{\,a}_i=f_{\theta}(\tilde{\mathbf{x}}^{\,h}_i).
\label{eq:cn_clean_adv_embedding}
\end{equation}
描述符网络的总损失定义为
\begin{equation}
\mathcal{L}_{F}
=
\mathcal{L}_{\mathrm{clean}}
+
\lambda_{\mathrm{adv}}\mathcal{L}_{\mathrm{adv}}
+
\lambda_{\mathrm{cons}}\mathcal{L}_{\mathrm{cons}},
\label{eq:cn_descriptor_loss}
\end{equation}
其中，干净样本识别损失和对抗样本识别损失分别为
\begin{align}
\mathcal{L}_{\mathrm{clean}}
&=
\mathcal{L}_{\mathrm{place}}
\left(
\left\{
f_{\theta}(\mathbf{x}_i)
\right\}_{i \in \mathcal{B}}
\right),
\label{eq:cn_clean_loss}
\\
\mathcal{L}_{\mathrm{adv}}
&=
\mathcal{L}_{\mathrm{place}}
\left(
\left\{
f_{\theta}\!\left(\tilde{\mathbf{x}}^{\,h}_i\right)
\right\}_{i \in \mathcal{B}}
\right).
\label{eq:cn_adv_loss}
\end{align}
为进一步约束同一样本在干净条件与遮挡条件下的全局表示不发生过大漂移，引入余弦一致性损失：
\begin{equation}
\mathcal{L}_{\mathrm{cons}}
=
\frac{1}{B}
\sum_{i=1}^{B}
\left(
1-
\frac{
\left\langle \mathbf{z}^{\,c}_i,\mathbf{z}^{\,a}_i \right\rangle
}{
\left\lVert \mathbf{z}^{\,c}_i \right\rVert_2
\left\lVert \mathbf{z}^{\,a}_i \right\rVert_2
}
\right).
\label{eq:cn_consistency_loss}
\end{equation}
式~\eqref{eq:cn_descriptor_loss} 中，$\mathcal{L}_{\mathrm{clean}}$ 用于维持网络在原始点云上的正常位置识别能力，
$\mathcal{L}_{\mathrm{adv}}$ 用于提升网络对结构化遮挡扰动的鲁棒性，而 $\mathcal{L}_{\mathrm{cons}}$ 则进一步增强同一场景在 clean/adv 两种观测条件下的表示一致性。

\paragraph{5) 交替优化过程}
本文采用逐批次交替优化策略。在每个训练迭代中，依次执行以下两个子步骤：
\begin{align}
\phi &\leftarrow \phi - \eta_{G} \nabla_{\phi}\mathcal{L}_{G},
\label{eq:cn_update_generator}
\\
\theta &\leftarrow \theta - \eta_{F} \nabla_{\theta}\mathcal{L}_{F},
\label{eq:cn_update_descriptor}
\end{align}
其中，$\eta_{G}$ 和 $\eta_{F}$ 分别表示生成器与描述符网络的学习率。

从训练机制上看，该交替优化过程对应一个稳定化的对抗博弈：首先，在固定描述符网络的条件下，生成器持续搜索更具破坏性的结构化动态遮挡；随后，在固定生成器的条件下，描述符网络通过干净样本监督、对抗样本监督以及表示一致性约束，逐步学习对该类遮挡的鲁棒表示。相较于同时更新两个网络，交替优化能够更清晰地分离“攻击者”和“防御者”的优化目标，从而提高训练稳定性并增强最终模型的泛化鲁棒性。
